% Generated by scripts/build_reports.py; edit reviewed Markdown inputs.
% Compile with XeLaTeX.
% Options for packages loaded elsewhere
\PassOptionsToPackage{unicode}{hyperref}
\PassOptionsToPackage{hyphens}{url}
\PassOptionsToPackage{dvipsnames,svgnames,x11names}{xcolor}
\documentclass[
  11pt,
]{article}
\usepackage{xcolor}
\usepackage[margin=1in]{geometry}
\usepackage{amsmath,amssymb}
\setcounter{secnumdepth}{-\maxdimen} % remove section numbering
\usepackage{iftex}
\ifPDFTeX
  \usepackage[T1]{fontenc}
  \usepackage[utf8]{inputenc}
  \usepackage{textcomp} % provide euro and other symbols
\else % if luatex or xetex
  \usepackage{unicode-math} % this also loads fontspec
  \defaultfontfeatures{Scale=MatchLowercase}
  \defaultfontfeatures[\rmfamily]{Ligatures=TeX,Scale=1}
\fi
\usepackage{lmodern}
\ifPDFTeX\else
  % xetex/luatex font selection
\fi
% Use upquote if available, for straight quotes in verbatim environments
\IfFileExists{upquote.sty}{\usepackage{upquote}}{}
\IfFileExists{microtype.sty}{% use microtype if available
  \usepackage[]{microtype}
  \UseMicrotypeSet[protrusion]{basicmath} % disable protrusion for tt fonts
}{}
\makeatletter
\@ifundefined{KOMAClassName}{% if non-KOMA class
  \IfFileExists{parskip.sty}{%
    \usepackage{parskip}
  }{% else
    \setlength{\parindent}{0pt}
    \setlength{\parskip}{6pt plus 2pt minus 1pt}}
}{% if KOMA class
  \KOMAoptions{parskip=half}}
\makeatother
\usepackage{longtable,booktabs,array}
\usepackage{caption}
\captionsetup[table]{skip=6pt}
\newcounter{none} % for unnumbered tables
\usepackage{calc} % for calculating minipage widths
% Correct order of tables after \paragraph or \subparagraph
\usepackage{etoolbox}
\makeatletter
\patchcmd\longtable{\par}{\if@noskipsec\mbox{}\fi\par}{}{}
\makeatother
% Allow footnotes in longtable head/foot
\IfFileExists{footnotehyper.sty}{\usepackage{footnotehyper}}{\usepackage{footnote}}
\makesavenoteenv{longtable}
\setlength{\emergencystretch}{3em} % prevent overfull lines
\providecommand{\tightlist}{%
  \setlength{\itemsep}{0pt}\setlength{\parskip}{0pt}}
\usepackage{amsmath,amssymb}
\usepackage{xurl}
\usepackage{fancyhdr}
\pagestyle{fancy}
\fancyhf{}
\fancyhead[L]{\small Pure Thought | Revised research specifications}
\fancyfoot[L]{\small 2026-09-17}
\fancyfoot[R]{\small\thepage}
\setlength{\headheight}{15pt}
\setlength{\emergencystretch}{3em}
\allowdisplaybreaks
\usepackage{bookmark}
\IfFileExists{xurl.sty}{\usepackage{xurl}}{} % add URL line breaks if available
\urlstyle{same}
% fallback for those not using the hyperref driver hyperxmp:
\makeatletter
\@ifundefined{xmpquote}{\newcommand{\xmpquote}[1]{#1}}{}
\makeatother
\hypersetup{
  pdftitle={28. Explicit finite-time action confinement from a normal form},
  pdfauthor={Pure Thought research portfolio},
  colorlinks=true,
  linkcolor={blue},
  filecolor={Maroon},
  citecolor={Blue},
  urlcolor={blue},
  pdfcreator={LaTeX via pandoc}}

\title{28. Explicit finite-time action confinement from a normal form}
\author{Pure Thought research portfolio}
\date{Revised 17 September 2026}

\begin{document}
\maketitle
\thispagestyle{empty}
\clearpage

\section{28. Explicit finite-time action confinement from a normal form}

\textbf{Revised specification --- 2026-09-11.}
\href{https://github.com/rick-stevens-ai/Pure-Thought/blob/main/reviews/28-Nekhoroshev-Stability-Theory.md}{Detailed
critique} \ensuremath{\cdot}
\href{https://github.com/rick-stevens-ai/Pure-Thought/blob/main/archive/original-PRDs/28-Nekhoroshev-Stability-Theory.md}{Original
description} \ensuremath{\cdot}
\href{https://github.com/rick-stevens-ai/Pure-Thought/blob/main/reviews/EVIDENCE-STANDARD.md}{Evidence
standard}

\subsection{Plain-language guide}

\subsubsection{The problem in everyday terms}

A system need not move on a perfectly orderly orbit to remain near its
initial state for a very long time. Small disturbances can produce only
a slow drift in certain quantities. This problem asks for an explicit
guarantee: how far can those quantities change, and for how long is that
bound valid? It focuses on a finite time horizon rather than proving
that a special motion persists forever.

\subsubsection{Key terms}

\begin{itemize}
\item
  \textbf{Near-integrable Hamiltonian} --- An energy function consisting
  of a well-understood integrable part plus a small perturbation.
\item
  \textbf{Action} --- A coordinate that is constant in the integrable
  reference system and may drift when a perturbation is added.
\item
  \textbf{Confinement bound} --- A guaranteed maximum change of the
  specified actions over a stated time interval.
\item
  \textbf{Normal form} --- A transformed expression of the equations
  separating manageable terms from a remainder.
\item
  \textbf{Canonical transformation} --- A change of variables preserving
  the Hamiltonian structure of the dynamics.
\item
  \textbf{Remainder} --- Terms not absorbed into the simplified normal
  form; bounding their size helps bound drift.
\item
  \textbf{Resonance} --- An integer relation among frequencies that can
  change how perturbations affect the motion.
\item
  \textbf{Nekhoroshev theorem} --- A theorem giving very long
  action-confinement times for small perturbations under specified
  regularity and nondegeneracy assumptions.
\item
  \textbf{Quasi-convexity} --- A particular curvature condition on the
  integrable energy function used by some such theorems.
\item
  \textbf{Dimensional time} --- A time expressed in declared physical or
  model units, rather than an unexplained numerical factor.
\end{itemize}

\subsubsection{Why this matters}

Finite-time guarantees answer a question that simulations alone cannot
settle beyond their run length. They also distinguish a valid stability
argument from an impressive-looking exponential formula with missing
constants or units. Understanding slow drift is relevant to the
mathematical study of planetary and other nearly integrable systems.

\subsubsection{What progress would mean}

A useful first result could be a modest but fully justified confinement
time. Improving that time with controlled remainders would be real
progress; asserting billions of years would require an appropriate
physical model and every conversion and bound needed to support it.

\subsection{Core question}

For a specified analytic Hamiltonian and action domain, what rigorously
computable action excursion and confinement time follow from a certified
normal form or quantitative Nekhoroshev theorem?

\subsection{Scope and assumptions}

Use H=h(I)+\ensuremath{\epsilon}f(I,\ensuremath{\theta}),
h=\ensuremath{\omega}\ensuremath{\cdot}I+|I|\ensuremath{{}^{2}}/2, in
two or three degrees of freedom on a compact action box with an explicit
complex extension. Supply \ensuremath{\omega}, a finite Fourier
polynomial f, and rational \ensuremath{\epsilon}. Bound the gradient
away from zero where the selected theorem requires it.

The initial target is a valid dimensional time bound, even if modest.
Planetary and post-Newtonian models are extensions requiring their own
reductions and hypotheses.

\subsection{Mathematical target}

For a quasi-convex theorem verify its actual condition, for example
v\ensuremath{{}^{T}}D\ensuremath{{}^{2}}h(I)v\ensuremath{\ge}m|v|\ensuremath{{}^{2}}
for v\ensuremath{\perp}\ensuremath{\nabla}h(I), plus the theorem's
gradient, domain and analyticity conditions. A positive Hessian
determinant alone is insufficient.

Alternatively construct a canonical \ensuremath{\Phi} with
H\ensuremath{\circ}\ensuremath{\Phi}=Z(J)+R(J,\ensuremath{\theta}),
bound |I\ensuremath{-}J|\ensuremath{\le}\ensuremath{\eta} and
||\ensuremath{\partial}\_\ensuremath{\theta}R||\ensuremath{\le}r on a
validated domain. Then
|I(t)\ensuremath{-}I(0)|\ensuremath{\le}2\ensuremath{\eta}+r|t| while
the transformed orbit remains in that domain. If resonant Z depends on
angles, include its action drift or restrict the claim to
conserved/protected components.

\subsection{Required outputs and proof obligations}

Export normal-form order, transformations, inverse/domain bounds,
Fourier tails, r, \ensuremath{\eta}, and a domain-containment proof for
the claimed interval. A Nekhoroshev result must include all constants
and its smallness threshold in T=C exp(c
\ensuremath{\epsilon}\textasciicircum{}\ensuremath{-}a) with stated
units; do not choose a from a desired time target.

\subsection{Validation and rejection controls}

Verify integrable \ensuremath{\epsilon}=0 behavior and a low-order
symbolic normal form. Check symplecticity with validated remainder
control. Confirm the action bound stays inside the certified domain.
Independently recompute
exp(10\textasciicircum{}0.3)\ensuremath{\approx}7.35: it is a
dimensionless factor, not 10\ensuremath{{}^{13}} years.

\subsection{Milestones and research extension}

Baseline: one explicit finite-time confinement certificate. Strong:
optimized truncation with a stronger bound or full quantitative
quasi-convex theorem application. Research extension: exponential
scaling across a parameter family; do not infer actual solar-system
stability from \ensuremath{\epsilon} alone.

Milestones are acceptance gates, not calendar promises. Before a run,
freeze the concrete instance/domain, tolerance, compute budget, and
permitted dependencies. A bounded search that fails to find a result
must preserve an unresolved outcome.

\subsection{Starting instruction}

Work on the baseline above. First make every input and convention
explicit, then validate the reference case. Build the checker alongside
the calculation. Label numerical evidence, finite-model certificates,
and continuum theorems separately; return the strongest justified result
and identify the exact unmet proof obligation.

\subsection{Primary references and their role}

\begin{itemize}
\tightlist
\item
  \href{https://arxiv.org/abs/1701.06026}{Zhang and Zhang, Improved
  stability for analytic quasi-convex nearly integrable systems} ---
  Theorem-specific quasi-convex exponential stability results.
\end{itemize}

References support definitions or limitations; the scope choices and
proposed milestones are editorial recommendations, not claims that these
results have already been obtained.

\subsection{Displayed formulation and claim boundary}

\href{https://github.com/rick-stevens-ai/Pure-Thought/blob/main/PRDs/28-Nekhoroshev-Stability-Theory.md}{Canonical
specification}.

\[H\circ\Phi=Z(J)+R(J,\theta),\qquad
|I-J|\le\eta,\qquad\|\partial_\theta R\|\le r,\]

\[|I(t)-I(0)|\le2\eta+r|t|.\]

This estimate assumes the angle-independent normal form shown and a
proved domain-containment interval. Include resonant drift if \(Z\)
depends on angles. A quantitative Nekhoroshev theorem instead requires
all its constants, hypotheses and threshold in
\(T=C\exp(c\epsilon^{-a})\), with units supplied by \(C\). The number
\(\exp(10^{0.3})\approx7.35413\) has no unit of time.

\subsection{Detailed review and rationale}

\href{https://github.com/rick-stevens-ai/Pure-Thought/blob/main/PRDs/28-Nekhoroshev-Stability-Theory.md}{Revised
problem} \ensuremath{\cdot}
\href{https://github.com/rick-stevens-ai/Pure-Thought/blob/main/archive/original-PRDs/28-Nekhoroshev-Stability-Theory.md}{Preserved
original} \ensuremath{\cdot}
\href{https://github.com/rick-stevens-ai/Pure-Thought/blob/f854d4e12b55c1f5b5c2a964b75f2d23496d981e/PRDs/28-Nekhoroshev-Stability-Theory.md}{Pinned
upstream source}

Line references below refer to the preserved original at commit
\texttt{f854d4e12b55}, not the rewritten file. ``Blocking'' means the
original issue can invalidate the scientific target or claimed
certificate; ``major'' means it materially changes scope, inference, or
acceptance criteria. These are critiques of specifications, not results
of running the embedded research-code sketches.

\subsubsection{Detailed findings}

\begin{enumerate}
\def\labelenumi{\arabic{enumi}.}
\item
  \textbf{Blocking --- incorrect steepness (30--31, 78--88).} Steepness
  and quasi-convexity are not equivalent to a positive Hessian
  determinant or an auxiliary convex function. The rewrite uses a
  theorem-specific condition.
\item
  \textbf{Blocking --- fabricated timescale (24, 111, 854--861).} The
  numerical exponential shown is about 7.35, not 10\ensuremath{{}^{13}}
  years, and no physical time unit or constants were supplied. This
  invalidates the stated solar-system benchmark.
\item
  \textbf{Major --- exponents (106--109).} ``Super-steep'' and
  \ensuremath{\epsilon}-dependent exponents are not a substitute for a
  cited precise theorem. Exponents cannot be optimized merely to meet a
  desired age target.
\item
  \textbf{Major --- global versus local control.} A truncated Fourier
  series, short integration, or local Hessian check does not certify the
  full action domain and time interval. A direct normal-form remainder
  bound supplies a defensible first milestone and exposes exactly what
  stronger claims require.
\end{enumerate}

\subsubsection{What the revision preserves and changes}

The central motivation is retained, but the executable question is now:
\textbf{For a specified analytic Hamiltonian and action domain, what
rigorously computable action excursion and confinement time follow from
a certified normal form or quantitative Nekhoroshev theorem?} The
revised scope chooses a concrete starting model; those choices are
proposed research-design decisions, not facts inferred from the
original. The baseline can succeed without resolving the full open
research program.

The replacement also removes the original implementation sketches from
the active specification. They remain in the archive for comparison, but
have not been made executable or certified. New implementations should
derive their equations from the revised definitions rather than copy
unchecked prototypes.

\subsubsection{Acceptance criteria to use instead}

Baseline: one explicit finite-time confinement certificate. Strong:
optimized truncation with a stronger bound or full quantitative
quasi-convex theorem application. Research extension: exponential
scaling across a parameter family; do not infer actual solar-system
stability from \ensuremath{\epsilon} alone.

\subsubsection{Sources supporting the corrections}

\begin{itemize}
\tightlist
\item
  \href{https://arxiv.org/abs/1701.06026}{Zhang and Zhang, Improved
  stability for analytic quasi-convex nearly integrable systems} ---
  Theorem-specific quasi-convex exponential stability results.
\end{itemize}

\section{Evidence and execution standard}

This standard applies to the revised specifications. It is a proposed
contract for future research runs; this editorial revision does not
supply the scientific certificates themselves.

\subsection{What ``pure thought'' means here}

Use explicitly stated mathematical models and deductions without fitting
conclusions to experimental outcomes. Synthetic examples, exact model
inputs, analytic references, numerical search, and established software
are allowed. Physical constants, electronic-structure approximations,
constitutive laws, and empirical parameter sets must be identified when
used. A model-derived result is conditional on that model.

The original ban on legacy software conflicts with its proposed reliance
on NumPy, SciPy, CVXPY, nauty, and other libraries. Reimplementing every
optimizer and algebra system is not necessary for scientific
independence. Prefer a small independently checkable certificate and a
documented trusted base. A fresh implementation can be a project
objective, but must not be confused with a mathematical proof or novelty
claim.

\subsection{Four evidence levels}

{\def\LTcaptype{none} % do not increment counter
\begin{longtable}[]{@{}
  >{\raggedright\arraybackslash}p{(\linewidth - 4\tabcolsep) * \real{0.3333}}
  >{\raggedright\arraybackslash}p{(\linewidth - 4\tabcolsep) * \real{0.3333}}
  >{\raggedright\arraybackslash}p{(\linewidth - 4\tabcolsep) * \real{0.3333}}@{}}
\toprule\noalign{}
\begin{minipage}[b]{\linewidth}\raggedright
Label
\end{minipage} & \begin{minipage}[b]{\linewidth}\raggedright
What it establishes
\end{minipage} & \begin{minipage}[b]{\linewidth}\raggedright
What it does not establish
\end{minipage} \\
\midrule\noalign{}
\endhead
\bottomrule\noalign{}
\endlastfoot
Numerical evidence & A reproducible computation at stated precision,
sample size, and discretization & Exact truth, continuum coverage, or a
universal theorem \\
Finite-model certificate & An exact/validated statement about the
specified finite model or finite relaxation & An omitted tail,
thermodynamic limit, or physical realization \\
Conditional mathematical theorem & A proof over the stated domain with
every hypothesis and error controlled & Truth outside those assumptions
or a demonstrated physical system \\
Formal proof & A checked theorem in a named proof assistant, with
imported axioms/dependencies listed & Correct physical modeling merely
by virtue of formalization \\
\end{longtable}
}

Formalization is an optional extra deliverable unless the specific
project makes it mandatory. High precision, JSON output, solver status
and small residuals do not by themselves confer any of these proof
levels.

\subsection{Claim directions}

For minimization over a larger feasible set, the optimum is a lower
bound on the original minimum. A physical feasible point supplies an
upper bound. For maximization the directions reverse. Track feasibility
of the primal/dual witnesses separately from their objective values.

Infeasibility of an inner ansatz excludes only that ansatz.
Infeasibility of an outer relaxation can exclude the original problem if
the map into the relaxation is proved and all approximations preserve
necessity. Passing necessary tests means ``not excluded''; existence
needs a constructive sufficient argument.

A smallest observed example is not a minimality theorem. A complete
catalog must define the universe, equivalence relation, boundaries, and
proof that pruning misses nothing. A finite-size numerical trend does
not prove an asymptotic limit.

\subsection{Minimum run manifest}

Before calculation, record the problem ID and revision hash, exact
inputs and conventions, quantified domain, equivalence relation, target
observable, accepted evidence level, tolerances, and resource limits.
Record dependencies and precision/rounding mode. Where randomization is
used, preserve seeds, trial counts and raw sufficient statistics.
Resource limits are planning decisions, not evidence that a mathematical
solution cannot exist.

An output bundle should include:

\begin{itemize}
\tightlist
\item
  The statement being claimed, including all assumptions and
  quantifiers.
\item
  Construction or proof object, plus an independently runnable verifier
  and its trust assumptions.
\item
  Numerical/error ledger separating roundoff, discretization,
  truncation, optimization, sampling and modeling error where relevant.
\item
  Positive reference cases and negative controls that the verifier must
  reject.
\item
  Coverage: included domain, omitted domain, unresolved regions and
  termination reason.
\end{itemize}

These are requirements for future implementations, not a promise of
polynomial-size certificates for every problem.

\subsection{Checking specific certificate types}

\textbf{LP/SDP/SOS:} verify equality constraints, objective direction,
and cone feasibility using exact arithmetic or outward-rounded
intervals. Rationalizing a numerical solution can break PSD constraints.
Correct residuals with justified norm bounds; do not silently project
and reuse the old objective. Distinguish a proof of a polynomial
approximation from a proof for the original function.

\textbf{Interval numerics:} state function enclosures, domains and
outward rounding. A small root box needs existence and, if claimed,
uniqueness. Cover the complement for a global enumeration. Include
singularities and tails rather than discarding them numerically.

\textbf{SAT/SMT:} a solver's UNSAT label is not a certificate. Preserve
a proof supported by a checker and the proof that the encoding, symmetry
reduction and blocked solutions match the mathematical universe. DRAT
applies to an appropriate Boolean CNF proof, not automatically to every
real-arithmetic or SDP claim.

\textbf{Statistics:} specify the estimand, dependence assumptions,
confidence procedure and bias. Zero observed failures does not prove
zero failure probability. MCMC error bars based on estimated
autocorrelation are not generally rigorous mixing guarantees.

\subsection{Milestones and novelty}

Baseline means a working, independently validated result on a concrete
reference problem. Strong means a broader or sharper result with its
claimed evidence level. Research extension means a possible novel result
whose feasibility and significance need review. A negative theorem can
be valuable; an unsuccessful search can also be useful if its scope is
documented, but the two are not interchangeable.

Do not promise ``publication-quality'' in a fixed number of months.
Before claiming novelty, conduct a focused current literature comparison
against the exact theorem/model/metric. The sources in this revision
check key definitions and limitations; they are not an exhaustive
novelty audit across all 30 fields.

\end{document}
